% ECONOMICS_PROBLEM: {"id": "D91-538", "title": "Identifying choice procedures", "jel_code": "D91", "source_id": "OP-0274"}
% !TeX program = xelatex
\documentclass[11pt,a4paper]{article}
\usepackage[margin=23mm]{geometry}
\usepackage{fontspec}
\setmainfont{Latin Modern Roman}
\usepackage{amsmath,amssymb,mathtools}
\usepackage{xcolor,enumitem,booktabs,longtable,array,xurl,hyperref}
\definecolor{muted}{HTML}{53636C}
\hypersetup{colorlinks=true,urlcolor=blue,linkcolor=blue}
\setlength{\parindent}{0pt}
\setlength{\parskip}{5pt}
\newcommand{\R}{\mathbb R}
\newcommand{\E}{\mathbb E}
\newcommand{\N}{\mathbb N}
\newcommand{\ind}{\mathbf 1}
\newcommand{\Var}{\operatorname{Var}}
\newcommand{\Cov}{\operatorname{Cov}}
\newcommand{\supp}{\operatorname{supp}}
\DeclareMathOperator*{\argmin}{arg\,min}
\DeclareMathOperator*{\argmax}{arg\,max}
\newcommand{\entrymeta}[3]{{\sffamily\small\color{muted}\textbf{\texttt{#1}}\quad|\quad #2\par #3\par}}
\newcommand{\Problemref}[1]{\texttt{#1}}
\newcommand{\JELref}[2]{\texttt{#2} (JEL \texttt{#1})}
\begin{document}
\section*{Identifying choice procedures}
\textbf{Problem ID:} \texttt{D91-538}\quad \textbf{JEL:} \texttt{D91}\par
% BEGIN ECONOMICS STATEMENT
\label{problem:OP-0274}
{\sffamily\small\color{muted}Original subject group: Economic theory, equilibrium and social choice\par}
\label{definitions:problem:30007555}
\textbf{Audit classification:} Published research agenda. Verified the survey’s §7 process-data and identification directions. The minimax KL design is editorial and requires restrictions excluding unavoidable overlap.
\subsection*{Definitions and precise research target}
\textbf{Model and research target.} Take a finite alternative set \(X\), a collection of menus \(\mathcal B\), and a set \(\mathcal T\) of feasible experimental interventions on presentation, search costs, and decision time. Let \(\mathcal M_1,\ldots,\mathcal M_K\) be specified parametric classes of choice procedures, including limited attention, satisficing, and uncertain preferences. Each pair \((k,\theta)\), where \(\theta\in\Theta_k\), induces a joint law \(P_{k,\theta}^{t,B}\) of chosen alternative and process observations under intervention \(t\) and menu \(B\). Choices alone may leave these classes observationally overlapping. An experiment \(e\) is a probability distribution over \(\mathcal T\times\mathcal B\), and its discrimination criterion is
\[
 J(e)=\inf_{k\ne\ell,\theta\in\Theta_k,\eta\in\Theta_\ell}
 \sum_{t,B}e(t,B)\,D_{\mathrm{KL}}(P_{k,\theta}^{t,B}\Vert P_{\ell,\eta}^{t,B}),
\]
where \(D_{\mathrm{KL}}\) is relative entropy. Identify when feasible interventions yield positive separation, characterize unavoidable equivalence classes, and design experiments maximizing discrimination subject to a stated cost budget. Establish sample requirements for valid model confidence sets under overlap, rather than forcing a single winner. Process measures must have explicitly modeled measurement error, and interventions must not covertly change the economic alternatives.

\emph{Definitions and maintained assumptions (editorial unless a source definition is named).} Take finite intervention set \(\mathcal T\), finite menu set \(\mathcal B\), and a common measurable observation space for choice and process data. For each model class fix a compact parameter space \(\Theta_k\), its joint observation laws, and a measurement-error kernel. Relative entropy is \(D_{\rm KL}(P\Vert Q)=\int\log(dP/dQ)\,dP\) when \(P\ll Q\), and \(+\infty\) otherwise; use a common dominating law and bounded densities when a finite continuous design criterion is desired. The design simplex is \(e(t,B)\ge0\), \(\sum e(t,B)=1\), with \(\sum e(t,B)c(t,B)\le C\). Two parameters are observationally equivalent if their laws coincide at every feasible intervention-menu pair. If classes share such a pair of parameters, the original global infimum \(J(e)\) is zero for every design; this is an elementary implication, not an open theorem. Meaningful uniform separation must concern disjoint closed hypothesis subsets or parameters at least \(\eta>0\) from the observational-equivalence set. Report confidence sets containing all compatible models.
\subsection*{Variants and assumption changes}
\begin{enumerate}
\item \textbf{Separated hypotheses (editorial).} Restrict to two compact disjoint observable-law sets at prescribed distance \(\eta\). Find cost-optimal designs and sample-complexity bounds for uniformly valid discrimination.
\item \textbf{Overlap (editorial).} Retain shared laws and use local alternatives approaching their intersection. Characterize equivalence classes and the rate at which procedure-specific counterfactual predictions become distinguishable with process measurements.
\end{enumerate}
\subsection*{References and source audit}
\begin{enumerate}
\item §7 pp.53–54 in February 2023 manuscript. Supports the stated research area; exact displayed specification is editorial. DOI publication is2024; full inspected manuscript is February2023. \url{https://geoffroydeclippel.net/Working%20Papers%20PDFs/JEL.pdf}
\end{enumerate}
\begin{enumerate}\raggedright
\item\hypertarget{definitions:source:30007555:1}{} Geoffroy de Clippel; Kareen Rozen. 2024. \emph{Bounded Rationality in Choice Theory: A Survey}. Published JEL 62(3),995–1039 (2024); inspected February2023 author manuscript §§6.4–7, printed pp.51–54.
\url{https://geoffroydeclippel.net/Working%20Papers%20PDFs/JEL.pdf}.
DOI: \url{https://doi.org/10.1257/jel.20231592}.
\end{enumerate}

\subsection*{Common conventions: Source mathematical conventions}

These conventions apply to each entry unless explicitly replaced.

\textbf{Models and identification.} A primitive model \(m\) specifies all probability laws, feasible choices, information, equilibrium and measurement rules used by its target. The admissible class \(\mathcal M\) is fixed independently of outcomes. If \(P_O(m)\) is its observable probability law and \(\tau(m)\) the target, its identified set at observable law \(P\) is
\[
 \mathcal I_\tau(P)=\{\tau(m):m\in\mathcal M,\ P_O(m)=P\}.
\]
Point identification means that this set is a singleton. An empty set signals incompatibility of data and assumptions. Coordinate bounds need not describe the joint set. Bounds are sharp only if every retained target is attainable or arbitrarily approachable by compatible models. Finite bounds require explicit restrictions on unobserved quantities; unrestricted confounding, hidden wealth, extrapolation or arbitrary equilibrium selection can yield uninformative sets.

\textbf{Causal interventions.} A potential outcome \(Y_i(p)\) is the outcome under a complete policy regime \(p\), including timing, financing, information and market spillovers. The target population and units are fixed before treatment. With interference, use the complete assignment vector or a justified exposure mapping; individual no-interference notation alone cannot identify an economy-wide intervention. A difference of mean potential outcomes does not require the cross-world joint distribution. Conditioning on post-treatment migration, survival, enrollment or employment changes the estimand and needs separate justification.

\textbf{Statistical statements.} An estimator, confidence set, forecast or test specifies a sampling law, model class and loss/coverage criterion. Uniform \((1-\alpha)\)-coverage means \(\inf_{m\in\mathcal M}\Pr_m\{\tau(m)\in C_n\}\geq1-\alpha\), or an explicitly asymptotic version. The whole parameter space trivially has coverage, so informativeness requires a stated width, risk or power objective. A forecasting improvement is relative to a named benchmark, information set and evaluation population.

\textbf{Equilibrium and optimization.} An equilibrium comprises feasible plans, optimality, beliefs where needed, budget constraints and all market/resource clearing. The primitives or a stated benchmark define each correspondence \(\mathcal E(p;m)\). Nonemptiness is assumed or proved, never inferred just from compact policy sets. A selected-equilibrium target uses a declared selection rule; a robust target quantifies over all permitted equilibria. Use suprema or infima unless attainment is established. An \(\varepsilon\)-optimal policy is feasible and within \(\varepsilon\) of the finite optimum value. Report an empty feasible set separately.

\textbf{Welfare and accounting.} Normative utility, interpersonal normalization, social weights, numeraire, discounting and the use of public receipts are primitives. Behavioral utility need not coincide with the welfare criterion. Household welfare includes ownership income and rents once; adding profits again to compensating variation generally double-counts them. A monetary equivalent is defined by an explicit transfer and price convention, with monotonicity and range conditions. Average effects, mechanism contributions, transfers and welfare losses are different targets. Infinite sums require convergence and dynamic feasible plans require terminal or no-Ponzi conditions.

\textbf{Source versions.} A working-paper date, revision date and journal publication year are distinguished. A DOI for a journal version is not automatically the DOI of an earlier manuscript. Locators refer to the stated version and printed pages where specified. A metadata-only check does not verify a claim about a full-text conclusion. Every entry's source subsection narrows the provenance claim accordingly.
% END ECONOMICS STATEMENT
\end{document}
